% =====================================================================
% 00_resume.tex — résumé, abstract et mots-clés
% =====================================================================
\chapter*{Résumé}
\addcontentsline{toc}{chapter}{Résumé}
\markboth{Résumé}{Résumé}

Ce rapport présente la conception et la réalisation d'un système d'information bancaire
en bases de données relationnelle-objet sous Oracle Database 21c Express Edition.
Le projet couvre la gestion des clients, des conseillers, des agences, des comptes
(courant, épargne et à terme), des opérations et des signataires. Il met en œuvre
les mécanismes objets d'Oracle : types objet, héritage avec \code{compte\_type}
déclaré \code{NOT FINAL NOT INSTANTIABLE}, références \code{REF} avec intégrité
référentielle, collections \code{VARRAY} et \emph{nested tables}, méthodes membres
et statiques, polymorphisme, type récursif de parrainage et vue objet.

La démarche suit sept étapes : cadrage, conception UML, création du schéma,
insertion du jeu de données, manipulations, requêtes SQL, programmation PL/SQL
et fonctionnalités objet avancées. Le schéma final comporte 19 types métier,
11 tables logiques et une table de stockage de \emph{nested table}, 4 séquences,
6 triggers, 9 sous-programmes catalogués et une vue objet. Le script unique
\code{projet\_bdda\_banque.sql} régénère l'ensemble des parties A à E et a été
rejoué intégralement : aucun objet invalide n'est signalé et les seules erreurs
\code{ORA-} observées correspondent à des tests d'exception volontaires.

Les tests valident les règles de gestion : contrôle des droits, plafond journalier
de retrait, découvert autorisé, commission inter-agences, alertes d'échéance,
audit des taux et des découverts, calcul des intérêts, clôture de compte et
mutation d'un conseiller. Le rapport documente également les limites du modèle,
les choix de conception et une préparation complète à la soutenance orale.

\vspace{0.5cm}
\noindent\textbf{Mots-clés :} Oracle 21c, base de données relationnelle-objet,
types objet, héritage, \code{REF}, \code{VARRAY}, \emph{nested table}, PL/SQL,
triggers, vue objet, UML.

\clearpage

\begin{otherlanguage}{english}
\chapter*{Abstract}
\addcontentsline{toc}{chapter}{Abstract}
\markboth{Abstract}{Abstract}

This report presents the design and implementation of a banking information
system using Oracle Database 21c Express Edition and its object-relational
features. The project manages customers, advisors, branches, accounts (current,
savings and term deposits), operations and signatories. It relies on object
types, inheritance with a \code{NOT FINAL NOT INSTANTIABLE} account supertype,
\code{REF} references with referential integrity, \code{VARRAY} and nested-table
collections, member and static methods, polymorphism, a recursive sponsorship
type and an object view.

The methodology follows seven stages: scoping, UML design, schema creation,
test-data loading, data manipulation, SQL queries, PL/SQL programming and
advanced object features. The final schema contains 19 business types, 11 logical
tables plus one nested-table storage table, 4 sequences, 6 triggers, 9 stored
subprograms and one object view. The single script
\code{projet\_bdda\_banque.sql} regenerates parts A to E and has been replayed
from scratch: no invalid object remains, and the only \code{ORA-} errors are
intentional exception tests.

The tests validate access rights, the daily withdrawal limit, the authorised
overdraft, the inter-branch commission, maturity alerts, rate and overdraft
auditing, interest calculation, account closure and advisor transfers. The
report also discusses the model's limitations, the design rationale and a
complete defence preparation.

\vspace{0.5cm}
\noindent\textbf{Keywords:} Oracle 21c, object-relational database, object types,
inheritance, \code{REF}, \code{VARRAY}, nested table, PL/SQL, triggers, object
view, UML.
\end{otherlanguage}

\clearpage
